% ============================================================
%  第一章 三维欧氏空间中的曲线
%  结构沿袭苏步青、胡和生《微分几何》第一章
%  记号约定：位置向量用 \boldsymbol{}，单位向量（T/N/B/e）用
%            \mathbf{}，曲率/挠率等标量不加粗。
% ============================================================

\section{曲线 \ 曲线的切向量\ 弧长}
\begin{mydefinition}{曲线}
	设(O;xyz)是$\mathbb{E}^3$ 中的笛卡尔坐标系，
	\begin{equation}
		\label{eq:curve}
		\begin{cases}
			x=x(t)\\
			y=y(t)\\
			z=z(t)
		\end{cases}
	\end{equation}
	于是式~\eqref{eq:curve} 给出了从$t\in(a,b)$到$\mathbb{E}^3$中$(x,y,z)$的连续可微映射
	\begin{equation}
		\label{eq:curve-map}
		t \mapsto (x(t),y(t),z(t))
	\end{equation}
	在这个映射下，t被映射到点$\mathbf{P}(x(t),y(t),z(t)$,(a,b)的像集就构成了$\mathbb{E}^3$中的一条\textbf{连续可微的曲线}$\mathbf{C}$,简称\textbf{曲线}。我们把t称为曲线的参数，式~\eqref{eq:curve}就是曲线$C$的参数方程。常把式~~\eqref{eq:curve}写成向量形式：
	\begin{equation}
		\label{eq:curve-vector}
		\mathbf{r}=\mathbf{r}(t)=(x(t),y(t),z(t))
	\end{equation}
	从而把曲线上的点$\mathbf{P}$称作点$\mathbf{r}(t)$,简称t点或$\mathbf{P}(t)$点
\end{mydefinition}
\begin{mydefinition}{曲线的切向量}
	按照参数增加的方向可以确定曲线的正向。称向量
	\begin{equation}
		\label{eq:tangent-vector}
		\frac{d\mathbf{r}(t)}{dt}=(\frac{dx(t)}{dt},\frac{dy(t)}{dt},\frac{dz(t)}{dt})
	\end{equation}
	为曲线$\mathbf{C}$在t点的\textbf{切向量}。

	若在$t=t_0$处$\frac{d\mathbf{r}(t_0)}{dt}\neq \mathbf{0}$则称参数为$t_0$的点$\mathbf{P}(t_0)$为\textbf{正则点}，否则称为\textbf{奇点}。

	当曲线上所有点都是正则点时，称曲线为\textbf{正则曲线}。
\end{mydefinition}
例如曲线$\mathbf{r}(t)=(a\cos(t),a\sin(t),bt)$ 是一条圆柱螺旋线，其中$a$为圆柱半径，$b$决定螺距，其图像如图~\ref{fig:helix}所示。
\begin{figure}[htbp]
	\centering
	\begin{tikzpicture}[
		x={(0.87cm,0.5cm)},
		y={(-0.87cm,0.5cm)},
		z={(0cm,1cm)},
		line cap=round,
		line join=round
	]
		% 坐标轴
		\draw[-{Stealth}] (-2.3,0,0) -- (2.4,0,0) node[above left] {$x$};
		\draw[-{Stealth}] (0,-2.3,0) -- (0,2.4,0) node[above right] {$y$};
		\draw[-{Stealth}] (0,0,-0.3) -- (0,0,3.0) node[right] {$z$};
		% 圆柱投影辅助圆（底面）
		\draw[thin, gray!60, dashed] plot[domain=0:360, samples=100, smooth]
			({1.6*cos(\x)}, {1.6*sin(\x)}, 0);
		% 螺旋线
		\draw[very thick, main] plot[domain=0:4*pi, samples=200, smooth]
			({1.6*cos(deg(\x))}, {1.6*sin(deg(\x))}, {0.19*\x});
		% 半径 a
		\draw[dashed] (0,0,0) -- (1.6,0,0) node[midway, above, sloped] {$a$};
	\end{tikzpicture}
	\caption{圆柱螺旋线 $\mathbf{r}(t)=(a\cos t,\,a\sin t,\,bt)$}
	\label{fig:helix}
\end{figure}
通过计算可知，这条曲线是正则曲线，无奇点。

再例如曲线$\mathbf{r}(t)=(t^3,t^2,0)$ 是$xy$平面中的半立方抛物线（尖点曲线），在$t=0$处出现尖点，不是正则曲线，如图~\ref{fig:semicubical}所示。
\begin{figure}[htbp]
	\centering
	\begin{tikzpicture}[
		x={(0.87cm,0.5cm)},
		y={(-0.87cm,0.5cm)},
		z={(0cm,1cm)},
		line cap=round,
		line join=round
	]
		% 坐标轴
		\draw[-{Stealth}] (-4,0,0) -- (4,0,0) node[above left] {$x$};
		\draw[-{Stealth}] (0,-0.6,0) -- (0,3.0,0) node[above right] {$y$};
		\draw[-{Stealth}] (0,0,-0.3) -- (0,0,0.6) node[right] {$z$};
		% 半立方抛物线（尖点曲线）
		\draw[very thick, main] plot[domain=-1.55:1.55, samples=200, smooth]
			({(\x)^3}, {(\x)^2}, 0);
		% 尖点
		\fill (0,0,0) circle (1.5pt);
		\node[below left] at (0,0,0) {$t=0$（尖点）};
	\end{tikzpicture}
	\caption{半立方抛物线 $\mathbf{r}(t)=(t^3,\,t^2,\,0)$}
	\label{fig:semicubical}
\end{figure}
注意：$\frac{d\mathbf{r}}{dt}=(3t^2,2t,0)$，在$t=0$处为零向量，故原点（尖点）不是正则点。

如果我们想采用另一个参数$\bar{t} $来表示同一条曲线，为了保证参数$\bar{t} $与$t$的一一对应，我们需要满足以下条件：
\begin{equation}
	\frac{d\bar{t}}{dt}\neq0
\end{equation}
为了使增加方向同时对应曲线正向，则要求
\begin{equation}
	\frac{d\bar{t}}{dt}> 0
\end{equation}
然后有链式法则可以知道
\begin{equation}
	\frac{d\mathbf{r}}{dt}=\frac{d\mathbf{r}}{d\bar{t}}\frac{d\bar{t}}{dt}
\end{equation}
参数t的正则点也必然是$\bar{t} $的正则点。即正则点的与否与参数无关。

\begin{mydefinition}{弧长}
	对于正则曲线$\mathbf{r}=\mathbf{r}(t)$，称
	\begin{equation}
		\label{eq:arc-length}
		s(t)=\int_{t_0}^{t}|\frac{d\mathbf{r}}{dt}|dt
	\end{equation}
	为曲线从参数$t_0$到$t$的弧长。其中
	$|\frac{d\mathbf{r}}{dt}|=\sqrt{\left(\frac{dx(t)}{dt}\right)^2+\left(\frac{dy(t)}{dt}\right)^2+\left(\frac{dz(t)}{dt}\right)^2}$
\end{mydefinition}
显然，弧长$s$是$t$的可微函数，且有$$\frac{ds}{dt}=|\frac{d\mathbf{r(t)}}{dt}|$$

对正则曲线，因为$\frac{d\mathbf{r(t)}}{dt}\neq0$，所以$\frac{ds}{dt}>0$所以可以取$s$作为新参数。此时由于$$1=\frac{ds}{ds}=|\frac{d\mathbf{r(s)}}{ds}|$$,所以曲线的切向量$\frac{d\mathbf{r(s)}}{ds}$是单位向量。反之，当切向量为单位向量时，计算弧长被简化。由式~\eqref{eq:arc-length}可知：
$$s=s(t)=\int_{t_0}^{t}|\frac{d\mathbf{r}}{dt}|dt=\int_{t_0}^{t}dt=t-t_0$$

常用\textbf{"撇"}表示关于弧长参数$s$的导数，如：$$\mathbf{r}'(s)=\frac{d\mathbf{r}}{ds}\quad ,\quad \mathbf{r}''(s)=\frac{d^2\mathbf{r}}{ds^2}$$

下面有一个重要定理：
\begin{myTheorem}{曲率的定义定理}
	设曲线$C:\ \mathbf{r}=\mathbf{r}(s)$的每点有一个单位向量$\mathbf{a}(s)$，则
	\begin{equation}
		|\mathbf{a}'(s)| =\lim_{\Delta s \to 0} \left|\frac{\Delta \theta}{\Delta s}\right|
	\end{equation}
	其中$\Delta\theta$为向量$\mathbf{a}(s)$转到$\mathbf{a}(s+\Delta s)$所转过的角度。
\end{myTheorem}
\begin{proof}
	由导数定义，
	\begin{equation}
		|\mathbf{a}'(s)|=\lim_{\Delta s\to 0}\frac{|\mathbf{a}(s+\Delta s)-\mathbf{a}(s)|}{\Delta s}
		=\lim_{\Delta s\to 0}\frac{|\Delta\mathbf{a}|}{\Delta s}
	\end{equation}
	设$\Delta\theta$为向量$\mathbf{a}(s)$转到$\mathbf{a}(s+\Delta s)$所转过的角度（图~\ref{fig:dtheta}）。由于$\mathbf{a}(s)$与$\mathbf{a}(s+\Delta s)$都是单位向量，由余弦定理，
	\begin{figure}[htbp]
		\centering
		\begin{tikzpicture}[scale=1.8]
			% 坐标轴
			\draw[-{Stealth}] (-0.15,0) -- (3.7,0) node[below] {$x$};
			\draw[-{Stealth}] (0,-0.75) -- (0,1.15) node[left] {$y$};
			% 一般曲线 C
			\draw[very thick, main] plot[domain=0.2:3.3, samples=120, smooth]
				({\x}, {0.5*sin(\x r)+0.2});
			% 曲线上两点
			\coordinate (A) at (1.0,0.6207);
			\coordinate (B) at (2.5,0.4992);
			\fill (A) circle (1.2pt);
			\fill (B) circle (1.2pt);
			\node[fill=white, inner sep=1pt, below, yshift=-2pt] at (A) {$P(s)$};
			\node[fill=white, inner sep=1pt, below, yshift=-2pt] at (B) {$P(s+\Delta s)$};
			% 切向量 a(s)
			\draw[-{Stealth}, very thick, second] (A) -- ++({0.7*cos(15.1)},{0.7*sin(15.1)})
				node[fill=white, inner sep=1pt, above right] {$\mathbf{a}(s)$};
			% 切向量 a(s+Δs)
			\draw[-{Stealth}, very thick, blue] (B) -- ++({0.7*cos(-21.8)},{0.7*sin(-21.8)})
				node[fill=white, inner sep=1pt, below right] {$\mathbf{a}(s+\Delta s)$};
			% 夹角 Δθ（仅画弧线，避免与切向量重叠）
			\draw[gray, thick, -{Stealth}] (A) ++(15.1:0.35)
				arc[start angle=15.1, end angle=-21.8, radius=0.35];
			\node[gray, fill=white, inner sep=1pt] at ($(A)+(-3.35:0.55)$) {$\Delta\theta$};
		\end{tikzpicture}
		\caption{一般曲线上两点切向量的夹角$\Delta\theta$}
		\label{fig:dtheta}
	\end{figure}
	\begin{equation}
		|\Delta\mathbf{a}|^2=|\mathbf{a}(s+\Delta s)|^2+|\mathbf{a}(s)|^2-2|\mathbf{a}(s+\Delta s)||\mathbf{a}(s)|\cos\Delta\theta
		=2(1-\cos\Delta\theta)=4\sin^2\frac{\Delta\theta}{2}
	\end{equation}
	从而
	\begin{equation}
		|\Delta\mathbf{a}|=2\left|\sin\frac{\Delta\theta}{2}\right|
	\end{equation}
	代入导数定义式，并利用$\lim_{x\to 0}\frac{\sin x}{x}=1$，得
	\begin{equation}
		|\mathbf{a}'(s)|=\lim_{\Delta s\to 0}\frac{2\left|\sin\frac{\Delta\theta}{2}\right|}{\Delta s}
		=\lim_{\Delta s\to 0}\frac{\left|\sin\frac{\Delta\theta}{2}\right|}{\frac{\Delta\theta}{2}}\cdot\frac{|\Delta\theta|}{\Delta s}
		=\lim_{\Delta s\to 0}\left|\frac{\Delta\theta}{\Delta s}\right|
	\end{equation}
	命题得证。
\end{proof}



\section{主法向量与从法向量 \ 曲率与挠率}
对曲线$\mathbf{r}=\mathbf{r}(s)$，用$\mathbf{T}(s)$表示单位切向量，即$$\mathbf{T}(s)=\mathbf{r}'(s)$$
上节末定理表明，单位向量$\mathbf{a}(s)$的转动速度（即单位长度上方向转过的角度）等于$|\mathbf{a}'(s)|$。特别地，当$\mathbf{a}(s)$取为单位切向量$\mathbf{T}(s)=\mathbf{r}'(s)$时，它的转动快慢恰好刻画了曲线在一点附近"弯曲"的程度：转角越大，曲线弯得越厉害。

\begin{mydefinition}{曲率}
	设$C:\ \mathbf{r}=\mathbf{r}(s)$是以弧长$s$为参数的正则曲线，称
	\begin{equation}
		\label{eq:curvature}
		\kappa(s)=|\mathbf{T}'(s)|=|\mathbf{r}''(s)|
	\end{equation}
	为曲线$C$在$s$点的\textbf{曲率}。由定理可知它亦等于单位切向量的转角对弧长的变化率：
	\begin{equation}
		\kappa(s)=\lim_{\Delta s\to 0}\left|\frac{\Delta\theta}{\Delta s}\right|
	\end{equation}
	其中$\Delta\theta$为切向量$\mathbf{T}(s)$转到$\mathbf{T}(s+\Delta s)$所转过的角度。
\end{mydefinition}

\begin{mydefinition}{曲率半径}
	当曲率$\kappa(s)\neq 0$时，称其倒数
	\begin{equation}
		\label{eq:radius-of-curvature}
		\rho(s)=\frac{1}{\kappa(s)}
	\end{equation}
	为曲线$C$在$s$点的\textbf{曲率半径}。它表示该点处密切圆的半径：曲率越大，曲率半径越小，曲线弯曲越剧烈；当曲线退化为直线时，$\kappa=0$，曲率半径趋于无穷大。
\end{mydefinition}

例如直线$\mathbf{r}(s)=\mathbf{r}(0)+\mathbf{v}\,s$（$|\mathbf{v}|=1$，$s$为弧长），其切向量为常向量
$$\mathbf{T}(s)=\mathbf{r}'(s)=\mathbf{v},\qquad \mathbf{T}'(s)=\mathbf{0}$$
故直线的曲率恒为零，$\kappa\equiv 0$，曲率半径$\rho=\frac{1}{\kappa}$趋于无穷大，这与"直线是弯曲程度最低的曲线"一致。

例如圆的参数方程$\mathbf{r}(s)=(R\cos\frac{s}{R},\,R\sin\frac{s}{R})$（$s$为弧长）满足
$$\mathbf{T}(s)=\mathbf{r}'(s)=\left(-\sin\frac{s}{R},\,\cos\frac{s}{R}\right),\qquad \kappa(s)=|\mathbf{T}'(s)|=\frac{1}{R}$$
故圆在各点的曲率恒为$\frac{1}{R}$，曲率半径恒为$R$，这正说明圆处处以相同的程度弯曲。

上节定理中取单位向量$\mathbf{a}(s)$为曲线的单位切向量$\mathbf{T}(s)=\mathbf{r}'(s)$，即可进一步建立曲线在一点处的活动标架。设$C:\ \mathbf{r}=\mathbf{r}(s)$是以弧长$s$为参数的正则曲线，且曲率$\kappa(s)=|\mathbf{T}'(s)|\neq 0$。

\begin{mydefinition}{主法向量}
	当$\kappa(s)\neq 0$时，称
	\begin{equation}
		\label{eq:principal-normal}
		\mathbf{N}(s)=\frac{\mathbf{T}'(s)}{\kappa(s)}
	\end{equation}
	为曲线$C$在$s$点的\textbf{主法向量}。由$|\mathbf{T}(s)|=1$知$\mathbf{T}(s)\cdot\mathbf{T}'(s)=0$，故$\mathbf{N}(s)\perp\mathbf{T}(s)$；又由$\kappa(s)=|\mathbf{T}'(s)|$知$\mathbf{N}(s)$为单位向量，且指向曲线弯曲（$\mathbf{T}$转动）的方向，即曲线的凹侧。
\end{mydefinition}

\begin{mydefinition}{密切平面}
	由$\mathbf{T}(s)$与$\mathbf{N}(s)$张成的平面称为曲线$C$在$s$点的\textbf{密切平面}（osculating plane）。它是经过$s$点的所有平面中与曲线"最贴近"的平面；等价地，它可由$s$点附近曲线上三点当三点都趋于$s$点时所确定的极限平面来定义。
\end{mydefinition}

\begin{mydefinition}{从法向量}
	称
	\begin{equation}
		\label{eq:binormal}
		\mathbf{B}(s)=\mathbf{T}(s)\times\mathbf{N}(s)
	\end{equation}
	为曲线$C$在$s$点的\textbf{从法向量}（副法向量）。它是垂直于密切平面的单位向量；$\{\mathbf{T}(s),\mathbf{N}(s),\mathbf{B}(s)\}$构成曲线在该点的右手单位正交活动标架，称为\textbf{Frenet标架}。
\end{mydefinition}

\begin{mydefinition}{从切平面}
	由$\mathbf{T}(s)$与$\mathbf{B}(s)$张成的平面称为曲线$C$在$s$点的\textbf{从切平面}（rectifying plane）。
\end{mydefinition}

此外，由$\mathbf{N}(s)$与$\mathbf{B}(s)$张成的平面称为曲线在$s$点的\textbf{法平面}（normal plane），它垂直于切向量$\mathbf{T}(s)$。上述三个平面两两垂直，其两两交线恰为$\mathbf{T},\mathbf{N},\mathbf{B}$三条坐标轴，如图~\ref{fig:frenet-frame}所示。

\begin{figure}[htbp]
	\centering
	\begin{tikzpicture}[
		scale=1.6,
		x={(0.87cm,0.5cm)},
		y={(-0.87cm,0.5cm)},
		z={(0cm,1cm)},
		line cap=round,
		line join=round
	]
		% 密切平面（T、N 张成，z=0）
		\fill[red!6] (0,0,0) -- (2,0,0) -- (2,1.5,0) -- (0,1.5,0) -- cycle;
		\draw[thin, gray!60] (0,0,0) -- (2,0,0) -- (2,1.5,0) -- (0,1.5,0) -- cycle;
		% 从切平面（T、B 张成，y=0）
		\fill[blue!6] (0,0,0) -- (2,0,0) -- (2,0,2) -- (0,0,2) -- cycle;
		\draw[thin, gray!60] (0,0,0) -- (2,0,0) -- (2,0,2) -- (0,0,2) -- cycle;
		% 法平面（N、B 张成，x=0）
		\fill[green!6] (0,0,0) -- (0,1.5,0) -- (0,1.5,2) -- (0,0,2) -- cycle;
		\draw[thin, gray!60] (0,0,0) -- (0,1.5,0) -- (0,1.5,2) -- (0,0,2) -- cycle;
		% 曲线（位于密切平面内）
		\draw[very thick, main] plot[domain=0:1.9, samples=120, smooth]
			({\x}, {0.25*\x*\x}, 0);
		% Frenet 标架
		\draw[-{Stealth}, very thick, second] (0,0,0) -- (1.8,0,0)
			node[right, fill=white, inner sep=1pt] {$\mathbf{T}$};
		\draw[-{Stealth}, very thick, third] (0,0,0) -- (0,1.2,0)
			node[left, fill=white, inner sep=1pt] {$\mathbf{N}$};
		\draw[-{Stealth}, very thick, blue] (0,0,0) -- (0,0,1.6)
			node[above, fill=white, inner sep=1pt] {$\mathbf{B}$};
		% 平面标注
		\node[fill=white, inner sep=1pt] at (1.0,0.75,0) {密切平面};
		\node[fill=white, inner sep=1pt] at (1.0,0,1.5) {从切平面};
		\node[fill=white, inner sep=1pt] at (0,0.75,1.5) {法平面};
		\fill (0,0,0) circle (1.5pt);
		\node[below left, fill=white, inner sep=1pt] at (0,0,0) {$P(s)$};
	\end{tikzpicture}
	\caption{曲线在一点处的Frenet标架与三个基本平面}
	\label{fig:frenet-frame}
\end{figure}

\section{Frenet标架 \ Frenet公式}

由前节的三个单位向量$\mathbf{T}(s)$、$\mathbf{N}(s)$、$\mathbf{B}(s)$组成的标架$\{\mathbf{T},\mathbf{N},\mathbf{B}\}$称为曲线$C$在$s$点的\textbf{Frenet标架}（活动标架）。由于$\mathbf{B}=\mathbf{T}\times\mathbf{N}$且三向量两两垂直，它构成曲线在该点的右手单位正交标架。

\begin{mydefinition}{挠率}
	先求$\mathbf{B}'(s)$的方向。因$\mathbf{B}=\mathbf{T}\times\mathbf{N}$，由$\mathbf{T}'=\kappa\mathbf{N}$得
	\begin{equation}
		\mathbf{B}'(s)\cdot\mathbf{T}(s)
		=\bigl(\mathbf{T}'(s)\times\mathbf{N}(s)\bigr)\cdot\mathbf{T}(s)
		=\kappa(s)\bigl(\mathbf{N}(s)\times\mathbf{N}(s)\bigr)\cdot\mathbf{T}(s)=0,
		\label{eq:挠率:B垂直T}
	\end{equation}
	又$|\mathbf{B}|=1$，故$\mathbf{B}'\cdot\mathbf{B}=0$；于是$\mathbf{B}'(s)$垂直于$\mathbf{T}$和$\mathbf{B}$，必平行于$\mathbf{N}(s)$。因此存在实数$\tau(s)$使
	\begin{equation}
		\mathbf{B}'(s)=-\tau(s)\,\mathbf{N}(s)
		\label{eq:挠率:定义}
	\end{equation}
	称$\tau(s)$为曲线$C$在$s$点的\textbf{挠率}。它刻画曲线偏离密切平面（即偏离平面曲线）的程度。
\end{mydefinition}

\begin{theorem}[Frenet--Serret公式]
	设$C:\ \mathbf{r}=\mathbf{r}(s)$是以弧长$s$为参数的正则曲线，曲率$\kappa(s)>0$，挠率$\tau(s)$，则Frenet标架满足方程组
	\begin{equation}
		\mathbf{T}'=\kappa\mathbf{N},\qquad
		\mathbf{N}'=-\kappa\mathbf{T}+\tau\mathbf{B},\qquad
		\mathbf{B}'=-\tau\mathbf{N},
		\label{eq:Frenet公式:方程组}
	\end{equation}
	称为\textbf{Frenet--Serret公式}。其矩阵形式为
	\begin{equation}
		\begin{pmatrix}
			\mathbf{T}'\\ \mathbf{N}'\\ \mathbf{B}'
		\end{pmatrix}
		=
		\begin{pmatrix}
			0 & \kappa & 0\\
			-\kappa & 0 & \tau\\
			0 & -\tau & 0
		\end{pmatrix}
		\begin{pmatrix}
			\mathbf{T}\\ \mathbf{N}\\ \mathbf{B}
		\end{pmatrix},
		\label{eq:Frenet公式:矩阵形式}
	\end{equation}
	系数矩阵是反对称矩阵。
\end{theorem}

\begin{proof}
	第一式即主法向量的定义式~\eqref{eq:principal-normal}。为证第二式，计算$\mathbf{N}'(s)$在Frenet标架下的分量。由$\mathbf{N}\cdot\mathbf{T}=0$、$\mathbf{N}\cdot\mathbf{B}=0$分别求导：
	\begin{equation}
		\mathbf{N}'\cdot\mathbf{T}=-\mathbf{N}\cdot\mathbf{T}'=-\kappa,
		\qquad
		\mathbf{N}'\cdot\mathbf{B}=-\mathbf{N}\cdot\mathbf{B}'=\tau,
	\end{equation}
	又$\mathbf{N}'\cdot\mathbf{N}=0$。故$\mathbf{N}'$在$\{\mathbf{T},\mathbf{N},\mathbf{B}\}$下的分量为$(-\kappa,\,0,\,\tau)$，即$\mathbf{N}'=-\kappa\mathbf{T}+\tau\mathbf{B}$。第三式即挠率定义式~\eqref{eq:挠率:定义}。命题得证。
\end{proof}

挠率还可直接用位置向量的各阶导数表出。若曲线以弧长$s$为参数，由$\mathbf{T}'=\kappa\mathbf{N}$与式~\eqref{eq:Frenet公式:方程组}可得
\begin{equation}
	\mathbf{r}'=\mathbf{T},\qquad
	\mathbf{r}''=\kappa\mathbf{N},\qquad
	\mathbf{r}'''=\kappa'\mathbf{N}+\kappa(-\kappa\mathbf{T}+\tau\mathbf{B})
	=-\kappa^{2}\mathbf{T}+\kappa'\mathbf{N}+\kappa\tau\mathbf{B}
\end{equation}
因而混合积
\begin{equation}
	(\mathbf{r}',\mathbf{r}'',\mathbf{r}''')
	=\mathbf{r}'\cdot(\mathbf{r}''\times\mathbf{r}''')
	=\kappa^{2}\tau
	\label{eq:挠率:混合积}
\end{equation}
由此得到
\begin{equation}
	\kappa(s)=|\mathbf{r}''(s)|,\qquad
	\tau(s)=\frac{(\mathbf{r}',\mathbf{r}'',\mathbf{r}''')}{|\mathbf{r}''|^{2}},
	\label{eq:挠率:任意参数公式}
\end{equation}
对任意参数$t$，则挠率为
\begin{equation}
	\tau(t)=\frac{(\mathbf{r}'(t),\mathbf{r}''(t),\mathbf{r}'''(t))}{|\mathbf{r}'(t)\times\mathbf{r}''(t)|^{2}}.
	\label{eq:挠率:一般参数公式}
\end{equation}

\begin{example}{圆柱螺旋线的曲率与挠率}
	对螺旋线$\mathbf{r}(t)=(a\cos t,\,a\sin t,\,bt)$，令$c=\sqrt{a^{2}+b^{2}}$，弧长为$s=ct$，其弧长参数方程为
	\begin{equation}
		\mathbf{r}(s)=\left(a\cos\frac{s}{c},\,a\sin\frac{s}{c},\,\frac{b}{c}s\right)
	\end{equation}
	求导得
	\begin{equation}
		\mathbf{T}=\left(-\frac{a}{c}\sin\frac{s}{c},\,\frac{a}{c}\cos\frac{s}{c},\,\frac{b}{c}\right),\qquad
		\mathbf{T}'=-\frac{a}{c^{2}}\left(\cos\frac{s}{c},\,\sin\frac{s}{c},\,0\right)
	\end{equation}
	于是曲率与主法向量为
	\begin{equation}
		\kappa=|\mathbf{T}'|=\frac{a}{c^{2}}=\frac{a}{a^{2}+b^{2}},\qquad
		\mathbf{N}=\frac{\mathbf{T}'}{\kappa}=-\left(\cos\frac{s}{c},\,\sin\frac{s}{c},\,0\right)
	\end{equation}
	又由$\mathbf{B}=\mathbf{T}\times\mathbf{N}$得
	\begin{equation}
		\mathbf{B}=\left(\frac{b}{c}\sin\frac{s}{c},\,-\frac{b}{c}\cos\frac{s}{c},\,\frac{a}{c}\right),\qquad
		\mathbf{B}'=-\frac{b}{c^{2}}\left(\cos\frac{s}{c},\,\sin\frac{s}{c},\,0\right)=-\frac{b}{c^{2}}\mathbf{N}
	\end{equation}
	故挠率为
	\begin{equation}
		\tau=\frac{b}{c^{2}}=\frac{b}{a^{2}+b^{2}}
	\end{equation}
	特别地，$\kappa$与$\tau$之比$\frac{\kappa}{\tau}=\frac{a}{b}$为常数，且当$b=0$时$\tau=0$，螺旋线退化为圆（平面曲线）；当$a=0$时$\kappa=0$，退化为直线。
\end{example}

由$\tau=0$恒成立即$\mathbf{B}$为常向量，可推出曲线是平面曲线的充要条件：

\begin{proposition}{平面曲线的判别}
	正则曲线$C$落在某平面内的充要条件是其挠率恒为零，即$\tau(s)\equiv 0$。
\end{proposition}

\begin{proof}
	若曲线是平面曲线，则其所有密切平面都与该平面重合，法向量$\mathbf{B}$为常向量，故$\mathbf{B}'=\mathbf{0}$，$\tau=0$。反之若$\tau\equiv 0$，则$\mathbf{B}'=-\tau\mathbf{N}=\mathbf{0}$，$\mathbf{B}$为常向量$\mathbf{B}_{0}$。于是$\mathbf{r}(s)\cdot\mathbf{B}_{0}$的导数为
	\begin{equation}
		\frac{\mathrm{d}}{\mathrm{d}s}\bigl(\mathbf{r}(s)\cdot\mathbf{B}_{0}\bigr)
		=\mathbf{T}(s)\cdot\mathbf{B}_{0}=0
	\end{equation}
	故$\mathbf{r}(s)\cdot\mathbf{B}_{0}$为常数，即曲线上所有点落在与$\mathbf{B}_{0}$垂直的平面内。命题得证。
\end{proof}
